\makeatletter
\@ifundefined{keepfig}{\newenvironment{keepfig}{\par\edef\keepfig@pi{\the\parindent}\noindent
  \begin{minipage}{\linewidth}\parindent=\keepfig@pi\relax}{\end{minipage}\par}}{}
\makeatother
\section{Многочлены с группой Галуа \texorpdfstring{$S_p$}{S\_p} над \texorpdfstring{$\Q$}{Q}}

Если $f\in\Q[X]$ неприводим степени $n$, то $G_f$~--- транзитивная подгруппа $S_n$ (\rf{3.27d}), и <<в типичном случае>> это
вся $S_n$. Но доказать, что у данного многочлена $G_f=S_n$, не так просто: нужно предъявить в $G_f$ достаточно много
перестановок. Для простой степени $p$ хватает двух: цикла длины $p$ и транспозиции. Цикл длины $p$ в $G_f$ есть всегда,
когда $f$ неприводим, а транспозицию даёт комплексное сопряжение, если у $f$ ровно два невещественных корня.

\begin{lemma}[4.15]
Пусть $p$ простое. Группа $S_p$ порождается любой транспозицией и любым циклом длины $p$.
\end{lemma}

\begin{proof}
Пусть $\tau$~--- транспозиция, $\sigma$~--- цикл длины $p$, и $H=\langle\sigma,\tau\rangle$. Перенумеровав элементы, можно
считать, что $\tau=(1\,2)$. Запишем цикл так, чтобы он начинался с $1$: $\sigma=(1\,i_2\,\dots\,i_p)$. Тогда
$\sigma^k(1)=i_{k+1}$, и при некотором $k$, $1\le k\le p-1$, имеем $\sigma^k(1)=2$. Перестановка $\sigma^k$ имеет порядок $p$
(так как $p$ простое и $p\nmid k$), а каждый элемент порядка $p$ в $S_p$~--- цикл длины $p$: порядок перестановки равен
наименьшему общему кратному длин её независимых циклов, поэтому все длины равны $1$ или $p$, а сумма длин равна $p$.
Значит, $\sigma^k=(1\,2\,j_3\,\dots\,j_p)$~--- снова цикл длины $p$, и он лежит в $H$. Заменим $\sigma$ на $\sigma^k$ и
перенумеруем элементы $j_3,\dots,j_p$ числами $3,\dots,p$; теперь $\tau=(1\,2)$, $\sigma=(1\,2\,3\,\dots\,p)$.

Для любой перестановки $\rho$ имеем $\rho\,(a\,b)\,\rho^{-1}=(\rho(a)\,\rho(b))$. Поэтому
\[
\sigma^i(1\,2)\sigma^{-i}=(i+1\;\;i+2)\in H,\qquad i=0,1,\dots,p-2,
\]
т.\,е.\ $H$ содержит все <<соседние>> транспозиции $(1\,2),(2\,3),\dots,(p-1\;p)$. Они порождают $S_p$: по индукции
$(a\;\,b+1)=(b\;\,b+1)(a\,b)(b\;\,b+1)$ для $a<b$, так что $H$ содержит все транспозиции, а каждая перестановка~---
произведение транспозиций.
\end{proof}

\begin{center}
\begin{keepfig}
\centering
\begin{tikzpicture}[every node/.style={font=\scriptsize}]
\foreach \s/\lab in {0/{$(1\,2)$},1/{$\sigma(1\,2)\sigma^{-1}=(2\,3)$},2/{$\sigma^2(1\,2)\sigma^{-2}=(3\,4)$}}{
\begin{scope}[xshift=\s*4.6cm]
\draw[gray] (0,0) circle (1.2);
\foreach \k in {1,...,7}{
  \coordinate (P\k) at ({90-(\k-1)*360/7}:1.2);
  \fill (P\k) circle (1.5pt);
  \node at ({90-(\k-1)*360/7}:1.5) {$\k$};}
\pgfmathtruncatemacro{\pa}{\s+1}\pgfmathtruncatemacro{\pb}{\s+2}
\draw[cnew,very thick,<->] (P\pa) -- (P\pb);
\draw[cold,->] (-30:0.55) arc (-30:-120:0.55);
\node[cold] at (0,0) {$\sigma$};
\node at (0,-2) {\lab};
\end{scope}}
\end{tikzpicture}
\figcap{Доказательство \rf{4.15} при $p=7$: цикл $\sigma=(1\,2\,\dots\,7)$~--- поворот на один шаг; сопряжение с $\sigma$ сдвигает
соседнюю транспозицию $(1\,2)$ по кругу, и получаются все соседние транспозиции, а они порождают $S_p$.}
\end{keepfig}
\end{center}

\begin{remark}[4.15a]
Простота $p$ существенна. В $S_4$ транспозиция $(1\,3)$ и цикл $(1\,2\,3\,4)$ порождают группу диэдра $D_4$ порядка $8$:
если расставить $1,2,3,4$ в вершинах квадрата по кругу, то $(1\,2\,3\,4)$~--- поворот, а $(1\,3)$~--- отражение относительно
диагонали $\{2,4\}$. Здесь степень $(1\,2\,3\,4)^2=(1\,3)(2\,4)$ цикла уже не является циклом длины $4$, и первый шаг
доказательства (перевести $1$ в $3$ степенью цикла, оставаясь в классе циклов длины $4$) не проходит: $\sigma^2(1)=3$ для
$\sigma=(1\,2\,3\,4)$, но $\sigma^2$~--- не цикл.
(С соседней транспозицией $(1\,2)$ цикл $(1\,2\,3\,4)$ порождает всю $S_4$~--- вторая половина доказательства годится при любом $n$.)
\end{remark}

\begin{proposition}[4.16]
Пусть $f\in\Q[X]$~--- неприводимый многочлен простой степени $p$, у которого ровно два невещественных корня в $\CC$.
Тогда $G_f=S_p$.
\end{proposition}

\begin{proof}
Пусть $E\subset\CC$~--- поле разложения $f$, порождённое над $\Q$ корнями $f$ в $\CC$, и $\alpha$~--- корень $f$. Так как $f$
неприводим, $[\Q[\alpha]:\Q]=p$, и по мультипликативности степеней (\rf{1.20}) $p$ делит $[E:\Q]=|G_f|$ (\rf{3.10}: $E/\Q$~---
расширение Галуа, так как $f$ сепарабелен в характеристике~$0$). По теореме Коши в $G_f$ есть элемент порядка $p$, т.\,е.\
(как в доказательстве \rf{4.15}) цикл длины $p$ на корнях.

По \rf{4.3} комплексное сопряжение задаёт элемент $G_f$, который оставляет на месте вещественные корни и меняет местами
два невещественных (они комплексно сопряжены друг другу). Значит, это транспозиция. По \rf{4.15} $G_f=S_p$.
\end{proof}

Условие <<ровно два невещественных корня>> использовалось только для того, чтобы найти в $G_f$ транспозицию. Поэтому
рассуждение доказывает больше.

\begin{corollary}[4.16a]
\textup{(a)} Транзитивная подгруппа $G\subset S_p$ \textup($p$ простое\textup), содержащая транспозицию, совпадает с $S_p$.

\textup{(b)} Если в транзитивной подгруппе $G\subset S_p$ есть перестановка, независимые циклы которой имеют длины
$2,\ell_1,\dots,\ell_r$ с нечётными $\ell_i$, то $G=S_p$.
\end{corollary}

\begin{proof}
(a) Орбита транзитивной группы состоит из $p$ элементов, так что $p$ делит $|G|$, в $G$ есть элемент порядка $p$, т.\,е.\
цикл длины $p$; остаётся применить \rf{4.15}. (b) Если $L$~--- наименьшее общее кратное чисел $\ell_i$ (оно нечётно), то в
$L$-й степени перестановки циклы длины $\ell_i$ исчезают, а цикл длины $2$ остаётся (так как $L$ нечётно); получается транспозиция.
\end{proof}

Вместе с теоремой Дедекинда (\rf{4.28}) пункт (b) даёт удобный признак: если неприводимый унитарный $f\in\Z[X]$ простой
степени $p$ по модулю некоторого простого $\ell$, не делящего дискриминант, раскладывается на неприводимый квадратный
множитель и множители нечётных степеней, то $G_f=S_p$. Пример~--- в \rf{4.18}.

Пример, к которому применимо \rf{4.16},~--- многочлен $X^5-6X+3$ (\rf{5.35e}). Построим такие многочлены для каждого простого
$p$. Нам понадобится наблюдение из анализа: при малом вертикальном сдвиге графика число его пересечений с осью абсцисс
не меняется.

\begin{lemma}[4.16b]
Пусть $g\in\R[X]$~--- многочлен степени $\ge1$, все вещественные корни которого простые, и $c_1<\dots<c_k$~--- все
вещественные корни $g'$. Положим $e=\min_j|g(c_j)|$ \textup($e=+\infty$, если $k=0$\textup). Тогда $e>0$, и при любом
вещественном $t$ с $|t|<e$ многочлен $g-t$ имеет столько же вещественных корней, сколько $g$, и все они простые.
\end{lemma}

\begin{proof}
Корни $g$ простые, т.\,е.\ не являются корнями $g'$ (\rf{2.19}); поэтому $g(c_j)\ne0$ и $e>0$. Точки $c_j$ делят $\R$ на
промежутки $I_0=(-\infty,c_1]$, $I_j=[c_j,c_{j+1}]$, $I_k=[c_k,+\infty)$ (при $k=0$~--- один промежуток $\R$); на каждом из
них $g'$ не меняет знака и обращается в нуль лишь в концах, так что $g$ строго монотонна. Пусть $|t|<e$. Решения
уравнения $g(x)=t$ не совпадают ни с одним $c_j$ (там $|g|\ge e>|t|$), поэтому каждое лежит внутри ровно одного
промежутка, внутри промежутка их не больше одного, и все они~--- простые корни $g-t$ (производная $g'$ в них не равна
нулю). Остаётся заметить, что попадание $t$ в образ $g(I_j)$ не зависит от $t\in(-e,e)$. Образ ограниченного
промежутка~--- отрезок с концами $g(c_j)$, $g(c_{j+1})$; оба конца по модулю $\ge e$, и отрезок содержит $t$ тогда и только
тогда, когда концы разных знаков. Образ луча $I_0$~--- луч с вершиной $g(c_1)$, идущий к $-\infty$ или к $+\infty$; он
содержит $t$ тогда и только тогда, когда луч идёт <<в сторону нуля>>, т.\,е.\ $g(c_1)>0$ при луче $(-\infty,g(c_1)]$ и
$g(c_1)<0$ при луче $[g(c_1),+\infty)$. Так же для $I_k$. При $k=0$ функция $g$ строго монотонна на $\R$, её степень
нечётна (у многочлена чётной степени $\ge2$ производная имеет нечётную степень и вещественный корень), и $g(x)=t$
имеет ровно одно решение при любом $t$. Во всех случаях число решений одно и то же для $t=0$ и для любого $t\in(-e,e)$.
\end{proof}

\begin{example}[4.17]
Пусть $p\ge5$ простое. Выберем чётное число $m>0$ и чётные целые числа $n_1<n_2<\dots<n_{p-2}$, и пусть
\[
g(X)=(X^2+m)(X-n_1)(X-n_2)\cdots(X-n_{p-2}).
\]
Вещественные корни $g$~--- это $n_1,\dots,n_{p-2}$ (так как $x^2+m>0$), и все они простые. По \rf{4.16b}
$e=\min_{g'(c)=0,\ c\in\R}|g(c)|>0$; выберем \emph{нечётное} натуральное $n$ с $\frac2n<e$ и положим
\[
f(X)=g(X)-\frac2n .
\]
\emph{Вещественные корни.} По \rf{4.16b} у $f$ ровно $p-2$ вещественных корня. Остальные два корня невещественны (и
комплексно сопряжены): $f$ раскладывается над $\R$ на $p-2$ линейных множителя и квадратный трёхчлен без вещественных
корней, корни которого лежат в $\CC$.

\emph{Неприводимость.} Рассмотрим многочлен $nf(X)=ng(X)-2=nX^p+a_1X^{p-1}+\dots+a_p\in\Z[X]$. По модулю $2$ имеем
$X^2+m\equiv X^2$, $X-n_i\equiv X$, т.\,е.\ $g\equiv X^p\pmod 2$: все коэффициенты $g$, кроме старшего, чётны. Поэтому
$a_1,\dots,a_{p-1}$ чётны, а старший коэффициент $n$ нечётен. Свободный член
\[
a_p=n\,g(0)-2=-n\,m\,n_1n_2\cdots n_{p-2}-2
\]
($p-2$ нечётно) чётен, но не делится на $4$: произведение $m\,n_1\cdots n_{p-2}$ содержит не меньше двух чётных
сомножителей и делится на $4$. По признаку Эйзенштейна (\rf{1.16}) для простого $2$ многочлен $nf$, а значит, и $f$,
неприводим над $\Q$. По \rf{4.16} $G_f=S_p$.

Каждое условие используется: $m>0$~--- чтобы у $X^2+m$ не было вещественных корней; чётность $m$ и $n_i$~--- чтобы
$g\equiv X^p\pmod2$ и $4\mid m\,n_1\cdots n_{p-2}$; нечётность $n$~--- чтобы старший коэффициент $nf$ не делился на $2$;
$\frac2n<e$~--- чтобы сдвиг сохранил число вещественных корней.
\end{example}

\begin{center}
\begin{keepfig}
\centering
\begin{tikzpicture}[x=1.6cm,y=0.13cm,every node/.style={font=\scriptsize}]
\draw[->] (-2.5,0) -- (2.6,0) node[right] {$x$};
\draw[->] (0,-19) -- (0,19);
\draw[gray,dashed] (-2.5,11.31) -- (2.5,11.31) node[right] {$e$};
\draw[gray,dashed] (-2.5,-11.31) -- (2.5,-11.31) node[right] {$-e$};
\draw[cold,thick,domain=-2.27:2.27,samples=120,smooth] plot (\x,{\x*\x*\x*\x*\x-2*\x*\x*\x-8*\x});
\draw[cnew,thick,domain=-2.24:2.29,samples=120,smooth] plot (\x,{\x*\x*\x*\x*\x-2*\x*\x*\x-8*\x-2});
\foreach \r in {-2,0,2}{\fill[cold] (\r,0) circle (1.6pt);}
\foreach \r in {-1.9556,-0.2464,2.0394}{\fill[cnew] (\r,0) circle (1.6pt);}
\node[cold] at (-1.414,14) {$g$};
\node[cnew] at (1.414,-16.3) {$f=g-2$};
\node[below left] at (-2,-0.5) {$-2$}; \node[below right] at (2,-0.5) {$2$};
\end{tikzpicture}
\figcap{Пример \rf{4.17a}: $g=(X^2+2)(X^3-4X)$ (синий) пересекает ось в $-2,0,2$; критические значения $\pm8\sqrt2$ по модулю
больше $2$, поэтому сдвинутый график $f=g-2$ (красный) тоже пересекает ось ровно трижды, а два корня $f$ невещественны.}
\end{keepfig}
\end{center}

\begin{example}[4.17a]
Возьмём $p=5$, $m=2$, $(n_1,n_2,n_3)=(-2,0,2)$: $g=(X^2+2)(X^3-4X)=X^5-2X^3-8X$. Производная $5X^4-6X^2-8=(5X^2+4)(X^2-2)$
имеет вещественные корни $\pm\sqrt2$, и $g(\pm\sqrt2)=\mp8\sqrt2$, так что $e=8\sqrt2>2$ и можно взять $n=1$:
\[
f=X^5-2X^3-8X-2,\qquad G_f=S_5 .
\]
Действительно, $f$~--- многочлен Эйзенштейна для $2$, и у него ровно три вещественных корня ($\approx-1{,}956$, $-0{,}246$,
$2{,}039$).
\end{example}

Выбор $n$ в \rf{4.17} можно не делать, если взять $m$ большим: тогда годится $n=1$, и доказательство даже проще.

\begin{proposition}[4.17b]
Пусть $p\ge3$ простое, $n_1<\dots<n_{p-2}$~--- чётные целые числа, $m$~--- чётное число с $m\ge2$ и
$2m>n_1^2+\dots+n_{p-2}^2$. Тогда у многочлена $f=(X^2+m)(X-n_1)\cdots(X-n_{p-2})-2$ ровно два невещественных корня,
и $G_f=S_p$.
\end{proposition}

\begin{proof}
Неприводимость~--- признак Эйзенштейна для $2$, как в \rf{4.17} (теперь $n=1$). Пусть $h=\prod(X-n_i)$, $g=(X^2+m)h$, $f=g-2$.

\emph{Не меньше $p-2$ вещественных корней.} В каждом нечётном целом $t$ все множители $t-n_i$~--- ненулевые нечётные
целые числа, так что $|h(t)|\ge1$; если при этом $h(t)>0$, то $g(t)\ge t^2+m\ge3$ и $f(t)>0$. Кроме того, $f(n_i)=-2<0$.
Так как $h$ имеет нечётное число $p-2$ простых корней и старший коэффициент $1$, он положителен на промежутках
$(n_1,n_2)$, $(n_3,n_4)$, \dots, $(n_{p-4},n_{p-3})$ и $(n_{p-2},+\infty)$. На каждом из $\frac{p-3}2$ ограниченных
промежутков $(n_j,n_{j+1})$ лежит нечётное число $t=n_j+1$ (так как $n_{j+1}\ge n_j+2$), и $f(n_j)<0<f(t)$, $f(t)>0>f(n_{j+1})$:
на промежутке не меньше двух корней. На $(n_{p-2},+\infty)$ есть ещё один корень, так как $f(n_{p-2})<0$, а
$f(x)\to+\infty$. Итого не меньше $2\cdot\frac{p-3}2+1=p-2$ различных вещественных корней.

\emph{Не все корни вещественны.} Пусть $\sigma_1=\sum n_i$, $\sigma_2=\sum_{i<j}n_in_j$. Тогда $h=X^{p-2}-\sigma_1X^{p-3}+\sigma_2X^{p-4}-\dots$ и
$g=X^p-\sigma_1X^{p-1}+(\sigma_2+m)X^{p-2}+\dots$; у $f$ те же три старших коэффициента (так как $p\ge3$). Для корней
$\alpha_1,\dots,\alpha_p$ многочлена $f$ формулы Виета дают
\[
\sum_k\alpha_k^2=\Bigl(\sum_k\alpha_k\Bigr)^2-2\sum_{k<l}\alpha_k\alpha_l=\sigma_1^2-2\sigma_2-2m=\sum_in_i^2-2m<0,
\]
что невозможно, если все $\alpha_k$ вещественны.

Невещественные корни идут комплексно сопряжёнными парами, поэтому число вещественных корней имеет ту же чётность, что
$p$; оно $\ge p-2$ и $<p$, т.\,е.\ равно $p-2$. По \rf{4.16} $G_f=S_p$.
\end{proof}

Например, при $(n_1,n_2,n_3)=(-2,0,2)$ годится $m=6$ (так как $12>8$): многочлен
$f=(X^2+6)(X^3-4X)-2=X^5+2X^3-24X-2$ имеет три вещественных корня ($\approx-1{,}974$, $-0{,}083$, $2{,}024$) и
$G_f=S_5$.

\begin{remark}[4.18]
Не следует думать, что многочлен с группой Галуа $S_p$ обязан иметь ровно два невещественных корня. Например, у
\[
f=X^5-5X^3+4X-1
\]
все корни вещественны, а $G_f=S_5$. Покажем это.

\emph{Все корни вещественны.} Здесь $X^5-5X^3+4X=(X+2)(X+1)X(X-1)(X-2)$, и $f$~--- этот многочлен, сдвинутый на $1$
вниз. Значения
\[
f(-2)=-1,\quad f(-\tfrac32)=\tfrac{73}{32},\quad f(-1)=-1,\quad f(0)=-1,\quad f(\tfrac12)=\tfrac{13}{32},\quad f(1)=-1,\quad f(3)=119
\]
дают пять перемен знака: на промежутках $(-2,-\frac32)$, $(-\frac32,-1)$, $(0,\frac12)$, $(\frac12,1)$, $(2,3)$
(при этом $f(2)=-1$). Значит, у $f$ пять вещественных корней ($\approx-1{,}954$, $-1{,}151$, $0{,}276$, $0{,}791$, $2{,}038$).
(Это частный случай \rf{4.16b}: критические значения многочлена $X^5-5X^3+4X$ по модулю не меньше $1{,}41>1$.)

\emph{Неприводимость.} По модулю $2$ имеем $\bar f=X^5+X^3+1$. У него нет корней в $\F_2$ ($\bar f(0)=\bar f(1)=1$), а
единственный неприводимый квадратный многочлен над $\F_2$~--- это $X^2+X+1$; по модулю него $X^3\equiv1$, и
$\bar f\equiv X^2+1+1=X^2\not\equiv0$. Значит, $\bar f$ неприводим над $\F_2$ (приводимый многочлен степени $5$ имеет
множитель степени $1$ или $2$), и по \rf{1.18} $f$ неприводим над $\Q$. Следовательно, в $G_f$ есть цикл длины $5$.

\emph{Транспозиция.} Анализ здесь ничего не даёт: поле разложения лежит в $\R$, комплексное сопряжение действует на нём
тождественно, а по \rf{3.30c} вещественность всех корней не противоречит и разрешимости. Нужна арифметика. Дискриминант
$f$ равен $38569$ (простое число; в частности, не квадрат, так что $G_f\not\subset A_5$). По модулю $19$
\[
f\equiv(X^2-2X+7)(X^3+2X^2+11X+8)\pmod{19},
\]
и оба множителя неприводимы над $\F_{19}$ (у квадратного дискриминант $-24\equiv14$~--- не квадрат по модулю $19$, у
кубического нет корней в $\F_{19}$). По теореме Дедекинда (\rf{4.28}) в $G_f$ есть перестановка корней с циклами длин $2$ и
$3$, и по \rf{4.16a}(b) $G_f=S_5$.
\end{remark}

\subsection*{Транзитивные подгруппы $S_5$}

Для многочлена степени $5$ группа Галуа~--- одна из транзитивных подгрупп $S_5$ (\rf{3.27d}). Перечислим их. Отождествим
множество из пяти элементов с $\F_5=\Z/5\Z$ и рассмотрим группу $\mathrm{AGL}_1(\F_5)$ аффинных преобразований
$x\mapsto ux+v$ ($u\in\F_5^\times$, $v\in\F_5$) из \rf{3.30a}. В ней лежат подгруппы
\[
C_5=\{x\mapsto x+v\},\qquad D_5=\{x\mapsto\pm x+v\},\qquad F_{20}=\mathrm{AGL}_1(\F_5)=\{x\mapsto ux+v\}
\]
порядков $5$, $10$, $20$. Группа $C_5$ порождена циклом $(0\,1\,2\,3\,4)$; $D_5$~--- группа симметрий правильного
пятиугольника с вершинами $0,1,2,3,4$ (по кругу); $F_{20}$ (\emph{группа Фробениуса} порядка $20$)~--- нормализатор
$C_5$ в $S_5$ (\rf{3.23a}).

\begin{proposition}[4.18a]
С точностью до сопряжённости транзитивные подгруппы $S_5$~--- это
\[
C_5\ (5),\qquad D_5\ (10),\qquad F_{20}\ (20),\qquad A_5\ (60),\qquad S_5\ (120)
\]
\textup(в скобках~--- порядки\textup). Первые три разрешимы, последние две~--- нет. Каждая транзитивная подгруппа $G$ содержит
силовскую $5$-подгруппу $P$ группы $S_5$, и $G$ разрешима тогда и только тогда, когда $G\subset N(P)$, где
$N(P)\cong F_{20}$~--- нормализатор $P$ в $S_5$.
\end{proposition}

\begin{proof}
Пусть $G\subset S_5$ транзитивна. Тогда $5$ делит $|G|$, и в $G$ есть подгруппа $P$ порядка $5$~--- силовская и в $S_5$
($120=5\cdot24$).

\emph{Разрешимые.} По \rf{3.30a} (условия (iv) и (v)) $G$ разрешима тогда и только тогда, когда при подходящей нумерации
$G\subset\mathrm{AGL}_1(\F_5)$ и $G\supset C_5$. Подгруппы между $C_5$ и $\mathrm{AGL}_1(\F_5)$ соответствуют подгруппам
факторгруппы $\mathrm{AGL}_1(\F_5)/C_5\cong\F_5^\times$ (гомоморфизм $(x\mapsto ux+v)\mapsto u$ из доказательства \rf{3.30a}), а
это циклическая группа порядка $4$ с подгруппами порядков $1,2,4$: $\{1\}$, $\{\pm1\}$, $\F_5^\times$. Получаем $C_5$, $D_5$,
$F_{20}$; они разрешимы. Перенумерация~--- это сопряжение в $S_5$.

\emph{Неразрешимые.} По \rf{3.30a} ((ii)$\Leftrightarrow$(v)) $|G|\ge25$, а $|G|$ делит $120$ и делится на $5$, так что
$|G|\in\{30,40,60,120\}$. В $A_5$ нет подгрупп индекса $k\le4$: действие $A_5$ левыми сдвигами на $k$ смежных классах
дало бы гомоморфизм $A_5\to S_k$ с ядром, отличным от $A_5$ (действие транзитивно, $k>1$); по простоте $A_5$ ядро
тривиально, что невозможно при $60>k!$. Подгруппа $G\cap A_5$ (ядро ограничения знака на $G$) имеет в $G$ индекс $1$ или
$2$. При $|G|=30$ или $40$ её порядок равен $15$, $20$, $30$ или $40$; порядок $40$ не делит $60$, а остальные дали бы
подгруппу индекса $4$, $3$ или $2$ в $A_5$~--- противоречие. При $|G|=60$ имеем $|G\cap A_5|\ge30$, значит, $G\cap A_5=A_5$ и $G=A_5$. При $|G|=120$ имеем $G=S_5$.
Группы $A_5$ и $S_5$ транзитивны и неразрешимы ($A_5$ проста и неабелева).

Наконец, $F_{20}=N(C_5)$: она нормализует $C_5$, а $|N(C_5)|=120/6=20$, так как в $S_5$ шесть силовских $5$-подгрупп
($24$ цикла длины $5$, по четыре в каждой подгруппе). Поэтому разрешимая транзитивная $G$ (после перенумерации
$G\subset\mathrm{AGL}_1(\F_5)$, $C_5\subset G$) содержится в нормализаторе своей силовской подгруппы $P=C_5$; обратно,
подгруппа $N(P)\cong F_{20}$ разрешима, а с ней и все её подгруппы.
\end{proof}

\begin{center}
\begin{keepfig}
\centering
\begin{tikzpicture}[every node/.style={font=\small}]
\node (C) at (0,0) {$C_5$}; \node (D) at (0,1.3) {$D_5$}; \node (F) at (-1.5,2.6) {$F_{20}$};
\node (A) at (1.5,2.6) {$A_5$}; \node (S) at (0,3.9) {$S_5$};
\draw (C) -- node[left] {$2$} (D); \draw (D) -- node[left] {$2$} (F); \draw (D) -- node[right] {$6$} (A);
\draw (F) -- node[left] {$6$} (S); \draw (A) -- node[right] {$2$} (S);
\draw[cnew,rounded corners,dashed] (-2.15,-0.4) rectangle (0.42,3.0);
\node[cnew,font=\footnotesize,align=center] at (-3.5,1.3) {разрешимые\\ $\subset N(P)$};
\end{tikzpicture}
\figcap{Транзитивные подгруппы $S_5$ (с точностью до сопряжённости) и индексы включений. Разрешимые (в рамке)~---
подгруппы группы Фробениуса $F_{20}$.}
\end{keepfig}
\end{center}

Для вычислений по теореме Дедекинда полезно знать, какие цикловые типы встречаются в каждой группе. Тип пишем
длинами независимых циклов: $2\,1^3$~--- транспозиция, $2^2\,1$~--- произведение двух независимых транспозиций и т.\,д.
Число элементов каждого циклового типа в каждой транзитивной подгруппе $S_5$ дано в таблице; цикловой тип при
сопряжении не меняется, так что таблица годится для любой подгруппы, сопряжённой с указанной.
\begin{center}
\small
\begin{tabular}{@{}lcccccccc@{}}
\toprule
 & $1^5$ & $2\,1^3$ & $2^2\,1$ & $3\,1^2$ & $3\,2$ & $4\,1$ & $5$ & порядок\\
\midrule
$C_5$ & $1$ & --- & --- & --- & --- & --- & $4$ & $5$\\
$D_5$ & $1$ & --- & $5$ & --- & --- & --- & $4$ & $10$\\
$F_{20}$ & $1$ & --- & $5$ & --- & --- & $10$ & $4$ & $20$\\
$A_5$ & $1$ & --- & $15$ & $20$ & --- & --- & $24$ & $60$\\
$S_5$ & $1$ & $10$ & $15$ & $20$ & $20$ & $30$ & $24$ & $120$\\
\bottomrule
\end{tabular}
\end{center}
Строка $S_5$~--- прямой подсчёт: транспозиций $\binom52=10$; произведений двух независимых транспозиций $\frac12\binom52\binom32=15$;
циклов длины $3$ всего $\binom53\cdot2=20$, и каждый вместе с транспозицией на двух оставшихся элементах даёт элемент
типа $3\,2$ (их тоже $20$); циклов длины $4$ $\binom54\cdot3!=30$; циклов длины $5$ $4!=24$. Группа $A_5$ состоит из чётных
перестановок, т.\,е.\ типов $1^5$, $2^2\,1$, $3\,1^2$, $5$. Для $\mathrm{AGL}_1(\F_5)$ всё видно сразу: сдвиг $x\mapsto x+v$, $v\ne0$,~---
цикл длины $5$; при $u\ne1$ преобразование $x\mapsto ux+v$ имеет ровно одну неподвижную точку $v/(1-u)$, а на остальных
четырёх точках его циклы имеют длину, равную порядку $u$ в $\F_5^\times$: $2$ при $u=-1$ (пять элементов, по числу $v$) и
$4$ при $u=2,3$ (десять элементов).

\begin{corollary}[4.18b]
Пусть $G\subset S_5$ транзитивна. Если $G$ содержит перестановку типа $2\,1^3$ или $3\,2$, то $G=S_5$; если типа
$3\,1^2$, то $G\supset A_5$; если типа $4\,1$, то $G=F_{20}$ или $S_5$. Группа $G$ разрешима тогда и только тогда, когда
в ней нет элементов порядка $3$.
\end{corollary}

\begin{proof}
Первые три утверждения читаются по таблице (с учётом \rf{4.18a}). Порядки $C_5$, $D_5$, $F_{20}$ не делятся на $3$, а $A_5$ и
$S_5$ содержат циклы длины $3$.
\end{proof}

\subsection*{Резольвенты}

Чтобы узнать, в какой из подгрупп лежит $G_f$, удобно иметь элементы поля разложения с заданным стабилизатором.
Пусть $f\in F[X]$~--- сепарабельный многочлен степени $n$, $x_1,\dots,x_n$~--- его корни в поле разложения $F_f$, и
$G=G_f$ отождествлена с подгруппой $S_n$ правилом $\tau(x_i)=x_{\tau(i)}$. Для $\pi\in S_n$ и $\Theta\in F[X_1,\dots,X_n]$
положим $(\pi\Theta)(X_1,\dots,X_n)=\Theta(X_{\pi(1)},\dots,X_{\pi(n)})$; это действие $S_n$
($\pi(\rho\Theta)=(\pi\rho)\Theta$), и для $\tau\in G$
\begin{equation}\label{eq:c4-res1}
\tau\bigl(\Theta(x_1,\dots,x_n)\bigr)=\Theta(x_{\tau(1)},\dots,x_{\tau(n)})=(\tau\Theta)(x_1,\dots,x_n).
\end{equation}

\begin{proposition}[4.18c: резольвента]
Пусть $H=\{\pi\in S_n\mid\pi\Theta=\Theta\}$~--- стабилизатор многочлена $\Theta$, $\pi_1H,\dots,\pi_kH$~--- левые смежные
классы \textup($k=(S_n:H)$\textup), $\theta_j=(\pi_j\Theta)(x_1,\dots,x_n)$ \textup(не зависит от выбора представителя\textup), и
\[
R(Y)=\prod_{j=1}^k(Y-\theta_j)\qquad\text{\textup(резольвента многочлена $\Theta$\textup)}.
\]
\begin{enumerate}[label=\textup{(\alph*)}]
\item $R\in F[Y]$.
\item Пусть $R$ сепарабелен. Тогда $G$ действует на корнях $R$ так же, как на $S_n/H$ левыми сдвигами; неприводимые
множители $R$ над $F$ отвечают $G$-орбитам на $S_n/H$; $\theta_j\in F$ тогда и только тогда, когда $G\subset\pi_jH\pi_j^{-1}$. В
частности, $R$ имеет корень в $F$ тогда и только тогда, когда $G$ содержится в подгруппе, сопряжённой с $H$.
\item Если поле $F$ бесконечно, то найдётся многочлен $u\in F[X]$ степени $<n$ такой, что элементы $y_i=u(x_i)$
попарно различны, многочлен $\tilde f=\prod(X-y_i)$ лежит в $F[X]$, имеет то же поле разложения и ту же группу
Галуа с тем же действием на корнях \textup($\tau(y_i)=y_{\tau(i)}$\textup), а резольвента $\Theta$ для $\tilde f$
сепарабельна.
\item \textup(Лагранж\textup) Пусть $R$ сепарабелен. Тогда стабилизатор $\theta_j$ в $G$ равен $G\cap\pi_jH\pi_j^{-1}$,
\[
F[\theta_j]=F_f^{\,G\cap\pi_jH\pi_j^{-1}},\qquad [F[\theta_j]:F]=(G:G\cap\pi_jH\pi_j^{-1}).
\]
В частности, каждый элемент $F_f$, неподвижный относительно $G\cap\pi_jH\pi_j^{-1}$, есть многочлен от $\theta_j$ с
коэффициентами из $F$.
\end{enumerate}
\end{proposition}

\begin{proof}
(a) По \eqref{eq:c4-res1} $\tau\in G$ переводит $\theta_j$ в $(\tau\pi_j\Theta)(x)=\theta_l$, где $\tau\pi_jH=\pi_lH$; левый сдвиг на
$\tau$ переставляет смежные классы, поэтому $\tau$ переставляет сомножители $R$, и коэффициенты $R$ неподвижны
относительно $G$, т.\,е.\ лежат в $F_f^G=F$ (\rf{3.17}).

(b) Если $R$ сепарабелен, то $\theta_j$ попарно различны, и биекция $\pi_jH\mapsto\theta_j$ согласована с действием $G$
(по доказательству (a)). По \rf{3.27d} орбиты $G$ на корнях $R$~--- множества корней неприводимых множителей. Далее,
$\theta_j\in F$ $\iff$ $\tau\theta_j=\theta_j$ для всех $\tau\in G$ $\iff$ $\tau\pi_jH=\pi_jH$ $\iff$ $\pi_j^{-1}\tau\pi_j\in H$ для
всех $\tau\in G$ $\iff$ $G\subset\pi_jH\pi_j^{-1}$.

(c) Пусть $t_0,\dots,t_{n-1}$~--- переменные и $Y_i=\sum_lt_lx_i^l$. Матрица $(x_i^l)$~--- матрица Вандермонда различных
$x_i$, она обратима, поэтому $(t_0,\dots,t_{n-1})\mapsto(Y_1,\dots,Y_n)$~--- обратимая линейная замена переменных над $F_f$. Для
$\pi_jH\ne\pi_lH$ многочлены $\pi_j\Theta\ne\pi_l\Theta$, поэтому $(\pi_j\Theta-\pi_l\Theta)(Y_1,\dots,Y_n)$~--- ненулевой многочлен
от $t$; ненулевые и $Y_i-Y_l$ при $i\ne l$. Их произведение~--- ненулевой многочлен от $t$ с коэффициентами в $F_f$, и
так как $F$ бесконечно, он не обращается в нуль в некоторой точке $t\in F^n$ (индукция по числу переменных: у
ненулевого многочлена от одной переменной конечное число корней). Положим $u=\sum t_lX^l$. Тогда $y_i=u(x_i)$ различны,
$\tau(y_i)=u(x_{\tau(i)})=y_{\tau(i)}$, коэффициенты $\tilde f$ неподвижны относительно $G$ и лежат в $F$, элемент $G$,
оставляющий все $y_i$ на месте, тождествен (тогда $\tau(i)=i$ для всех $i$), так что $F[y_1,\dots,y_n]=F_f$ (\rf{3.17}), а значения
$(\pi_j\Theta)(y_1,\dots,y_n)$ попарно различны.

(d) Как в (b), для $\tau\in G$ имеем $\tau\theta_j=\theta_j\iff\tau\pi_jH=\pi_jH\iff\tau\in\pi_jH\pi_j^{-1}$ (здесь нужна
попарная различность $\theta_l$). Поэтому $\Gal(F_f/F[\theta_j])=G\cap\pi_jH\pi_j^{-1}$, и по основной теореме (\rf{3.17})
$F[\theta_j]$~--- неподвижное поле этой подгруппы, а его степень над $F$ равна её индексу в $G$. Поле $F[\theta_j]$
состоит из многочленов от $\theta_j$, так как $\theta_j$ алгебраичен над $F$ (\rf{1.27c}).
\end{proof}

Знакомые примеры: $\Theta=\prod_{i<j}(X_i-X_j)$ со стабилизатором $A_n$ (при $\operatorname{char}F\ne2$) даёт $R=Y^2-D(f)$ и
критерий <<$G_f\subset A_n\iff D(f)$~--- квадрат>> (\rf{4.1}--\rf{4.5}); $\Theta=X_1X_2+X_3X_4$ со стабилизатором $D_4$ даёт
кубическую резольвенту многочлена степени $4$ (\rf{4.8a}, \rf{4.9}). Для многочленов степени $5$ нужна подгруппа $F_{20}$.

Почему для степени $5$ не ищут, как для степени $4$, сюръективный гомоморфизм из $S_5$ на меньшую симметрическую группу?
У $S_5$ всего три нормальные подгруппы: $\{1\}$, $A_5$, $S_5$. Действительно, если $N\ne\{1\}$ нормальна в $S_5$, то
$N\cap A_5$ нормальна в простой группе $A_5$. При $N\cap A_5=A_5$ имеем $N\supset A_5$, т.\,е.\ $N=A_5$ или $S_5$. При
$N\cap A_5=\{1\}$ знак инъективен на $N$, так что $|N|=2$; но нормальная подгруппа порядка $2$ лежит в центре (её
неединичный элемент переходит при сопряжении в себя), а центр $S_5$ тривиален. Поэтому, кроме изоморфизмов, у $S_5$
есть лишь один нетривиальный гомоморфный образ~--- $S_5/A_5\cong\{\pm1\}$ (знак), и ему отвечает дискриминант.
Сюръекция $S_4\to S_3$ с ядром $V$ (\rf{4.7j}(b))~--- особенность малой степени; правильная точка зрения~--- действие на
смежных классах $S_n/H$: группа $S_4$ переставляет три разбиения корней на пары, т.\,е.\ классы $S_4/D_4$, а $S_5$
переставляет шесть классов $S_5/F_{20}$ (\rf{4.18d}). Последнее действие транзитивно и точно: его ядро
$\bigcap_\pi\pi F_{20}\pi^{-1}$ нормально в $S_5$ и содержится в $F_{20}$, т.\,е.\ тривиально. Оно задаёт вложение
$S_5\hookrightarrow S_6$ с транзитивным образом (не сопряжённым со стабилизатором точки $S_5\subset S_6$). Так же, перебирая
подгруппы $H\subset S_n$ и проверяя, есть ли у соответствующих резольвент корни в $F$, вычисляют группы Галуа и в
больших степенях.

\begin{lemma}[4.18d]
Пусть
\[
U=(X_1X_2+X_2X_3+X_3X_4+X_4X_5+X_5X_1)-(X_1X_3+X_3X_5+X_5X_2+X_2X_4+X_4X_1)
\]
\textup(сумма по сторонам пятиугольника $1\,2\,3\,4\,5$ минус сумма по его диагоналям\textup) и $\Theta=U^2$. Если
$\operatorname{char}F\ne2$, то стабилизатор $\Theta$ в $S_5$~--- группа Фробениуса $F_{20}$ \textup(для нумерации $i\leftrightarrow i-1\in\F_5$\textup), и $\Theta$ имеет
$6$ различных образов $\pi\Theta$.
\end{lemma}

\begin{proof}
Симметрии пятиугольника (группа $D_5$) переводят стороны в стороны и диагонали в диагонали и сохраняют $U$.
Преобразование $x\mapsto2x$ множества $\F_5$ переводит сторону $\{x,x+1\}$ в диагональ $\{2x,2x+2\}$ и диагональ
$\{x,x+2\}$ в сторону $\{2x,2x+4\}$, т.\,е.\ $U\mapsto-U$ и $\Theta\mapsto\Theta$. Значит, стабилизатор $H\supset F_{20}$, и $H$
транзитивна. По \rf{4.18a} $|H|\in\{20,60,120\}$ (порядок делится на $20$). При $|H|=20$ имеем $H=F_{20}$. При $|H|=60$
группа $H$ сопряжена с $A_5$, т.\,е.\ $H=A_5$ (подгруппа $A_5$ нормальна), но $A_5\not\supset F_{20}$: в $F_{20}$ есть
нечётные перестановки~--- циклы длины $4$. Наконец, $H\ne S_5$: при $(X_1,\dots,X_5)=(0,1,2,1,0)$ имеем $U=(2+2)-1=3$, а
после транспозиции $(1\,2)$, т.\,е.\ при $(1,0,2,1,0)$, имеем $U=2-(2+1)=-1$, и $\Theta$ меняется: $9-1=8\ne0$, так как
$\operatorname{char}F\ne2$. (При $\operatorname{char}F=2$ многочлен $U=\sum_{i<j}X_iX_j$ симметричен, и лемма неверна.) Итак,
$H=F_{20}$, и образов $(S_5:F_{20})=6$.
\end{proof}

Резольвента $R$ этого $\Theta$ имеет степень $6$ и называется \emph{резольвентой Кэли}. Шести смежным классам $S_5/F_{20}$
отвечают шесть сопряжённых с $F_{20}$ подгрупп, т.\,е.\ нормализаторы шести силовских $5$-подгрупп $S_5$; на языке
пятиугольников~--- шесть способов разбить десять пар из $\{1,\dots,5\}$ на <<стороны>> и <<диагонали>> (пятиугольник и
дополнительный к нему пятиугольник дают одно разбиение).

\begin{theorem}[4.18e: Кэли]
Пусть $\operatorname{char}F\ne2$, $f\in F[X]$~--- неприводимый сепарабельный многочлен степени $5$ и его резольвента Кэли
$R$ сепарабельна. Тогда либо $R$ неприводим над $F$, и $G_f\in\{A_5,S_5\}$, либо $R$~--- произведение линейного и
неприводимого многочлена степени $5$, и $G_f$~--- одна из групп $C_5,D_5,F_{20}$. Вместе с дискриминантом это даёт:
\begin{center}
\small
\begin{tabular}{@{}lll@{}}
\toprule
& $D(f)$ не квадрат в $F$ & $D(f)$~--- квадрат в $F$\\
\midrule
$R$ не имеет корней в $F$ & $S_5$ & $A_5$\\
$R$ имеет корень в $F$ & $F_{20}$ & $D_5$ или $C_5$\\
\bottomrule
\end{tabular}
\end{center}
В частности, при $\operatorname{char}F=0$ уравнение $f=0$ разрешимо в радикалах тогда и только тогда, когда $R$ имеет
корень в $F$.
\end{theorem}

\begin{proof}
Группа $G=G_f$ транзитивна (\rf{3.27d}) и содержит силовскую подгруппу $P$ порядка $5$ (\rf{4.18a}). Сопоставим смежному классу
$\pi F_{20}$ силовскую подгруппу $\pi C_5\pi^{-1}$. Это корректно определённая биекция $S_5/F_{20}$ на множество шести
силовских $5$-подгрупп (так как $F_{20}=N(C_5)$: $\pi C_5\pi^{-1}=\rho C_5\rho^{-1}\iff\rho^{-1}\pi\in N(C_5)$), согласованная с
действием $G$: левый сдвиг на $\tau$ переходит в сопряжение с $\tau$. Нормализатор $N(Q)$ силовской подгруппы $Q$ имеет
порядок $20$, и $Q$~--- его единственная (нормальная) силовская $5$-подгруппа. Подгруппа $P$ оставляет на месте силовскую подгруппу $Q$ тогда и только тогда, когда $P\subset N(Q)$, т.\,е.\
$P=Q$. Значит, $P$ неподвижна на одной точке из шести, а остальные пять образуют одну $P$-орбиту (орбиты группы порядка $5$
имеют $1$ или $5$ элементов).

Если $G$ разрешима, то по \rf{4.18a} $G\subset N(P)$, так что $G$ оставляет точку $P$ на месте, а остальные пять точек
образуют одну орбиту (она содержит $P$-орбиту из пяти точек). Если $G=A_5$ или $S_5$, то $G$ содержит все $24$ цикла длины
$5$, т.\,е.\ все шесть силовских подгрупп, и по теореме Силова они сопряжены в $G$: действие транзитивно. По \rf{4.18c}(b)
получаем утверждение о разложении $R$. Строки таблицы различают $\{C_5,D_5,F_{20}\}$ и $\{A_5,S_5\}$, столбцы~--- лежит ли
$G$ в $A_5$ (\rf{4.1}--\rf{4.5}); $C_5,D_5\subset A_5$, а $F_{20}\not\subset A_5$ (таблица цикловых типов). Последнее утверждение~---
теорема \rf{5.34} вместе с \rf{4.18a}.
\end{proof}

Если $R$ имеет кратный корень, а поле $F$ бесконечно (например, $F=\Q$), то по \rf{4.18c}(c) многочлен $f$ сначала заменяют на $\tilde f$ с той же группой.
Группы $C_5$ и $D_5$ резольвента не различает: $G_f=C_5$ тогда и только тогда, когда $|G_f|=5$, т.\,е.\ $f$ раскладывается
на линейные множители над $F[x_1]$. При $F\subset\R$ полезно и такое наблюдение: если у $f$ есть невещественный корень, то
комплексное сопряжение~--- элемент порядка $2$ в $G_f$, и $G_f\ne C_5$.

\begin{proposition}[4.18f: резольвента Кэли трёхчлена]
Для $f=X^5+aX+b$
\[
R(Y)=Y^6-40aY^5+880a^2Y^4-8960a^3Y^3+44800a^4Y^2-(108544a^5+3200000b^4)\,Y+102400a^6 .
\]
\end{proposition}

\begin{proof}
Рассмотрим $\prod_j(Y-\pi_j\Theta)$ с переменными $X_1,\dots,X_5$: перестановки переменных переставляют сомножители
(как в доказательстве \rf{4.18c}(a)), поэтому коэффициенты~--- симметрические многочлены от $X_i$ с целыми коэффициентами, и
по \rf{5.36} они~--- многочлены с целыми коэффициентами от элементарных симметрических многочленов $e_1,\dots,e_5$. Для
$f=X^5+aX+b$ имеем $e_1=e_2=e_3=0$, $e_4=a$, $e_5=-b$; поэтому коэффициенты $R$~--- многочлены от $a,b$ с целыми
коэффициентами, одни и те же для всех $a,b$. Припишем $x_i$ вес $1$; тогда $a$ имеет вес $4$, $b$~--- вес $5$, $\Theta$~--- вес $4$, и коэффициент при $Y^{6-k}$
однороден веса $4k$. Мономы $a^ib^j$ веса $4k\le24$: $a^k$, а при $k=5,6$ ещё $b^4$ и $ab^4$. Остаётся найти восемь
числовых коэффициентов. Это конечное вычисление; мы сделали его на компьютере, вычислив с высокой точностью корни и
значения $\theta_j$ при восьми парах целых $(a,b)$, и проверили результат в независимой точке. Коэффициент при $ab^4$
оказывается нулевым.
\end{proof}

\begin{example}[4.18g]
Проверим таблицу \rf{4.18e} на трёхчленах (все приводимые ниже утверждения проверены вычислением; резольвенты
сепарабельны).
\begin{enumerate}[label=(\arabic*),leftmargin=*]
\item $X^5-5X+12$: неприводим по модулю $7$; $D=2^{12}5^6=8000^2$; $R=(Y-100)(Y^5+300Y^4+\dots)$. Группа $D_5$ или $C_5$, а так
как вещественный корень один, $G_f=D_5$ (ср.\ \rf{2.4d}, где это доказано иначе).
\item $X^5+15X+12$: многочлен Эйзенштейна для $3$; $D=2^{10}3^45^5$~--- не квадрат; $R=(Y-180)(Y^5-420Y^4+\dots)$.
Значит, $G_f=F_{20}$, и уравнение разрешимо в радикалах.
\item $X^5+20X+16$: по модулю $3$ это неприводимый многочлен $X^5+2X+1$; $D=2^{16}5^6=32000^2$; у $R$ нет корней даже по
модулю $7$, а рациональный корень унитарного $R\in\Z[Y]$ был бы целым (\rf{1.14}). Значит, $G_f=A_5$.
\item $X^5-X-1$: неприводим по модулю $5$ (\rf{2.22b}); $D=2869=19\cdot151$; у $R=Y^6+40Y^5+880Y^4+8960Y^3+44800Y^2-3091456Y+102400$
нет корней по модулю $7$. Значит, $G_f=S_5$ (ср.\ \rf{4.30}).
\item $X^5-2$: $a=0$, $R=Y(Y^5-51200000)$, $D=2^45^5$~--- не квадрат; $G_f=F_{20}$ (ср.\ \rf{3.23}, \rf{3.23a}).
\end{enumerate}
Для многочлена $X^5+X^4-4X^3-3X^2+3X+1$ с корнем $2\cos\frac{2\pi}{11}$ (не трёхчлена) резольвента, вычисленная по
определению, имеет корень $0$, $D=11^4$~--- квадрат, и действительно $G_f=C_5$: поле $\Q[\cos\frac{2\pi}{11}]$~--- подполе
степени $5$ абелева расширения $\Q[\zeta_{11}]$ и потому расширение Галуа. В примерах \rf{4.17a} и \rf{4.18} резольвенты Кэли
неприводимы, а дискриминанты не квадраты, в согласии с $G_f=S_5$.
\end{example}

\begin{aside}[4.18h: уравнение пятой степени и модулярные функции]
Многочлен с группой $A_5$ или $S_5$ не решается в радикалах, но решается, если к радикалам добавить одну
<<трансцендентную операцию>>~--- обращение эллиптической модулярной функции (Эрмит, Кронекер, Бриоски, 1858). Это
аналог тригонометрического решения кубического уравнения в неприводимом случае (\rf{4.7g}; ср.\ \rf{3.22f}), где извлечение кубического
корня заменяется делением угла на три: $\cos(\frac13\arccos t)$. Схема такова.
\begin{enumerate}[label=(\arabic*),leftmargin=*]
\item Преобразованием Чирнгауза $y=h(x)$ ($h$~--- многочлен; ср.\ \rf{4.7i}) с коэффициентами, находимыми квадратными и кубическими радикалами
(Бринг, Джеррард), уравнение пятой степени приводится к виду $y^5+by+c=0$; при $b\ne0$ замена $y=\lambda x$, $\lambda^4=-b$,
даёт $x^5-x-a=0$, $a=-c/\lambda^5$ (а при $b=0$ уравнение $y^5=-c$ решается в радикалах).
\item Пусть $q=e^{i\pi\tau}$, $\operatorname{Im}\tau>0$, и $u(\tau)=\sqrt2\,q^{1/8}\prod_{n\ge1}\frac{1+q^{2n}}{1+q^{2n-1}}$ (это корень
четвёртой степени из эллиптического модуля $k$). Шесть величин $v_\infty=-u(5\tau)$, $v_m=u\bigl(\frac{\tau+16m}5\bigr)$,
$m\in\F_5$,~--- корни <<модулярного уравнения пятого порядка>>; его группа Галуа над $\CC(u)$~--- $\mathrm{PSL}_2(\F_5)\cong A_5$ (Галуа),
действующая на $\mathbb P^1(\F_5)=\{\infty,0,\dots,4\}$. Как и для кубической резольвенты, из шести корней строятся пять
величин, переставляемых этой группой: $y_r=(v_\infty-v_r)(v_{r+1}-v_{r+4})(v_{r+2}-v_{r+3})$, $r\in\F_5$. Эрмит
обнаружил, что они~--- корни уравнения вида Бринга--Джеррарда
\[
y^5-2000\,u^4(1-u^8)^2\,y-1600\sqrt5\,u^3(1-u^8)^2(1+u^8)=0,\qquad u=u(\tau).
\]
\item По данному $a$ модуль $k$ находится квадратными радикалами: равенство
$a=\frac{2(1+k^2)}{5^{5/4}\sqrt k\,\sqrt{1-k^2}}$ после возведения в квадрат, деления на $k^2$ и подстановки $w=k-k^{-1}$
становится квадратным уравнением $4w^2+25\sqrt5\,a^2w+16=0$, а $k$ находится из $k^2-wk-1=0$. Затем $\tau=i\,K(k')/K(k)$, где $k'=\sqrt{1-k^2}$ и $K$~--- полный эллиптический интеграл
(обращение модулярной функции), и корни $x^5-x-a$ равны $y_r/s$, где $s^4=2000u^4(1-u^8)^2$ и $s$ выбран так, чтобы
$1600\sqrt5\,u^3(1-u^8)^2(1+u^8)=a s^5$.
\end{enumerate}
Например, для $x^5-x-1$ корень $w\approx-0{,}29233$ даёт $k\approx0{,}86446$, $\tau\approx0{,}78414\,i$, и формулы дают все
пять корней ($1{,}16730\dots$, $-0{,}76488\pm0{,}35247i$, $0{,}18123\pm1{,}08395i$); невязки меньше $10^{-38}$. Годится и
второй корень $w\approx-13{,}683$: он даёт $k\approx0{,}07270$, $\tau\approx2{,}55057\,i$ и те же пять корней (в другом
порядке), так что выбор корня $w$ не важен. Тождество Эрмита мы проверили численно с точностью до $50$ знаков в
нескольких точках $\tau$. Доказательство требует теории модулярных функций и выходит за рамки этого текста.
\end{aside}

\begin{aside}[4.18i: явное решение разрешимой квинтики]
В отличие от \rf{4.18h}, при разрешимой группе $G_f$ корни выражаются в радикалах, и их можно выписать явно. Спускаемся по ряду
подгрупп с циклическими факторами; на каждом шаге циклическое расширение степени $d$ (после присоединения корней из
единицы степени $d$) порождается \emph{резольвентой Лагранжа}, $d$-я степень которой лежит в нижнем поле (\rf{5.30l} и абзац
после него; при $d=3$ это \rf{4.7d}, \rf{4.7e}). Разберём пример $f=X^5+X^4-4X^3-3X^2+3X+1$, $G_f=C_5$ (\rf{4.18g}).

\emph{Корни и группа.} Пусть $\zeta=e^{2\pi i/11}$. Корни $f$~--- числа $2\cos\frac{2\pi a}{11}$, $a=1,\dots,5$ (сопряжённые с
$2\cos\frac{2\pi}{11}=\zeta+\zeta^{-1}$). Так как $2^5\equiv-1\pmod{11}$, их удобно занумеровать так:
\[
\eta_j=\zeta^{2^j}+\zeta^{-2^j}=2\cos\frac{2\pi\cdot2^j}{11},\qquad j\in\Z/5\Z
\]
($\eta_{j+5}=\eta_j$; при $j=0,\dots,4$ числа $\pm2^j$ пробегают все ненулевые вычеты по модулю $11$). Ограничение на
$\Q[\eta_0]$ автоморфизма $\zeta\mapsto\zeta^2$ поля $\Q[\zeta]$~--- образующая $\sigma$ группы $G_f$, и $\sigma(\eta_j)=\eta_{j+1}$.

\emph{Присоединяем корни из единицы.} Пусть $\varepsilon=e^{2\pi i/5}=\frac{\sqrt5-1}4+\frac i4\sqrt{10+2\sqrt5}$ и
$L=\Q[\varepsilon,\eta_0]$. Степень $[L:\Q]$ делится на $[\Q[\varepsilon]:\Q]=4$ и на $[\Q[\eta_0]:\Q]=5$ и не больше $20$, т.\,е.\
равна $20$. Поэтому $L$~--- поле разложения $f$ над $\Q[\varepsilon]$ степени $5$, и ограничение
$\Gal(L/\Q[\varepsilon])\to G_f$ инъективно, а значит, биективно: $\sigma$ продолжается до автоморфизма $L$, тождественного
на $\Q[\varepsilon]$ (обозначим его тоже $\sigma$). Точно так же $L$~--- поле разложения $\Phi_5$ над $\Q[\eta_0]$ степени $4$, и
для каждого $k\in(\Z/5\Z)^\times$ есть автоморфизм $\rho_k$ поля $L$, тождественный на $\Q[\eta_0]$, с
$\rho_k(\varepsilon)=\varepsilon^k$.

\emph{Резольвенты Лагранжа.} Положим
\[
r_k=\sum_{j\in\Z/5\Z}\varepsilon^{jk}\eta_j,\qquad k\in\Z/5\Z .
\]
Тогда $\sigma(r_k)=\sum_j\varepsilon^{jk}\eta_{j+1}=\varepsilon^{-k}r_k$ и $\rho_k(r_1)=r_k$. Значит, $r_0=\sum_j\eta_j=-1$, а
$r_1^5$ и $c_k=r_k/r_1^k$ неподвижны относительно $\sigma$ и по \rf{3.17} лежат в $\Q[\varepsilon]$. Здесь $r_1\ne0$: иначе и
$r_k=\rho_k(r_1)=0$ при $k\ne0$, а обратное преобразование
\[
\eta_j=\frac15\sum_{k\in\Z/5\Z}\varepsilon^{-jk}r_k
\]
(оно следует из $\sum_k\varepsilon^{mk}=0$ при $m\not\equiv0\pmod5$) дало бы $\eta_j=-\frac15$ для всех $j$. Вычисление даёт
\[
r_1^5=11\,(-26-20\varepsilon+15\varepsilon^2-10\varepsilon^3),\qquad
c_2=-\frac{2+2\varepsilon+3\varepsilon^2+4\varepsilon^3}{11},
\]
\[
c_3=-\frac{10+7\varepsilon-2\varepsilon^2+4\varepsilon^3}{121},\qquad
c_4=-\frac{6-20\varepsilon-10\varepsilon^2-35\varepsilon^3}{1331},
\]
и формула обращения при $j=0$ даёт ответ в радикалах:
\[
2\cos\frac{2\pi}{11}=\eta_0=\frac15\bigl(-1+r_1+c_2r_1^2+c_3r_1^3+c_4r_1^4\bigr),\qquad r_1=\sqrt[5]{11\,(-26-20\varepsilon+15\varepsilon^2-10\varepsilon^3)} .
\]
Корень пятой степени здесь нужно выбрать правильно. При замене $r_1$ на $\varepsilon^mr_1$ (другой корень пятой степени
из того же числа) каждое $r_k=c_kr_1^k$ умножается на $\varepsilon^{mk}$, и формула даёт
$\frac15\sum_k\varepsilon^{mk}r_k=\eta_{-m}$: пять выборов корня дают ровно пять корней $f$. Получилась башня
$\Q\subset\Q[\sqrt5]\subset\Q[\varepsilon]\subset\Q[\varepsilon,r_1]=L$ из двух квадратичных шагов и одного шага Куммера
степени $5$.

Числа $r_k$, $k\ne0$,~--- это \emph{суммы Гаусса} $\sum_{a\in(\Z/11\Z)^\times}\chi^k(a)\,\zeta^a$ для характера $\chi$ циклической
группы $(\Z/11\Z)^\times=\langle2\rangle$, заданного равенством $\chi(2^j)=\varepsilon^j$ (так как $\chi(-1)=\chi(2^5)=1$, слагаемые
с $a$ и $-a$ объединяются в $\varepsilon^{jk}\eta_j$); для них $|r_k|^2=11$. Для $G_f=A_5$ или $S_5$ такой спуск обрывается
сразу после присоединения $\sqrt{D(f)}$: простая группа $A_5$ не имеет нормальных подгрупп, кроме $\{1\}$ и $A_5$, а значит,
и следующего циклического шага (ср.\ \rf{4.18e}).
\end{aside}
